\section*{Resumen}
\addcontentsline{toc}{section}{Resumen}
Se estudia la clasificación binaria de ventas unitarias mensuales de Wide World Importers (WWI), a partir de una extracción de su Data Mart de Ventas. La respuesta indica si las unidades observadas de un producto superan la mediana de sus tres meses previos. Se auditaron 228\,265 líneas de factura y se construyeron observaciones producto-mes con información disponible antes del horizonte objetivo. El perceptrón simple y la regresión logística se implementaron mediante NumPy sobre las mismas 20 características transformadas, con particiones temporales de 4\,599 observaciones de entrenamiento, 2\,628 de validación y 1\,111 de prueba.

La selección utilizó únicamente validación; las predicciones finales se conservaron y su consistencia se comprobó sin reentrenar. En prueba, la regresión logística alcanzó balanced accuracy de 62,42\%, frente a 54,10\% del perceptrón, 50,00\% del baseline mayoritario y 51,31\% de la regla histórica. La logística registró 370 falsos positivos y 49 falsos negativos, lo que limita su utilidad para alertas automáticas. El perceptrón mantuvo actualizaciones hasta la última época; la logística mostró descenso estable de entropía cruzada, sin demostración de convergencia completa al agotar 200 épocas. PCA se empleó solo para visualización y conservó 49,62\% de la varianza en dos componentes.

Las conclusiones son descriptivas de cinco meses de prueba con observaciones dependientes. La cobertura al cierre y la identidad por nombre son supuestos o limitaciones pendientes de validación operacional. Este documento recoge la primera aplicación efectivamente implementada en Ventas; la extensión a los demás Data Marts es una obligación académica pendiente, no una actividad opcional ni un resultado ya obtenido.

\noindent\textbf{Palabras clave:} Inteligencia de Negocios; Data Warehouse; clasificación temporal; perceptrón; regresión logística; PCA; trazabilidad.
